\section{结论}
本文依次建立原始素材的特征时间组织、局部缺失预测与可解释性分析模型，三问通过统一的内容掩码与时间回映规则相互连接，形成从素材时间组织到不完整特征预测、再到决策依据复核的处理框架。

第一，建立词级时间锚点与交叠池化模型，完成附件1全部100条视频的自主特征提取与时序组织。1932个原词中1920个生成时间区间，区间生成率为99.38\%。三模态特征、观测掩码、父词索引与帧级记录共同支持全量结果追溯，典型样本给出了文本片段、语音时段与视频帧段的逐项对应。

第二，基于附件2的官方对齐特征建立教师--学生模型TSRF-D，完成局部缺失条件下的极性分类与强度回归，验证集Macro-F1为0.6012、强度MAE为0.5915。受控遮挡实验显示：文本缺口引起最大的预测退化，8锚点文本缺口使测试集F1由0.6138降至0.5198，而同等条件下语音、视觉缺口分别为0.6157与0.6100；连续缺口跨度与缺失数量共同增加时性能总体下降，在同一实际新增遮挡量下连续缺口的退化更为明显。标准化特征的重建改善主要来自视觉模态，各组成部分的作用通过缺失条件下的对照实验给出。附件3全部30条样本均已完成预测，输出负向7条、中性11条、正向12条，并保存逐样本概率与强度。

第三，采用模态子集训练与固定参考类别的输入干预构建ICF-Shapley模型，完成附件4全部20条样本的预测与解释，输出中性7条、负向7条、正向6条。96条验证样本中，高分片段累计删除引起的参考类别概率下降大于随机对照，AOPC差为0.1068；模态贡献在多种子与多基线下具有较高一致性。通过原词区间与视频PTS，局部证据可回映至原始素材，形成预测、模态作用与可回看证据的对应记录。

上述结果说明：以词为公共锚点的时序组织能够同时支持特征提取、缺失建模与证据回映三项目标，三态掩码表示使局部缺失无需插补值即可参与训练，固定参考类别的联盟分解与片段删除可在同一预测器上给出可量化、可复核的解释。三问共享的时间口径与掩码约定也为后续扩展到更长缺口、更多模态与多任务联合预测提供了可直接复用的接口。
